%-------------------------------%
%  Author: Alessandro Sciarra   %
%    Date: 03 Jun 2026          %
%-------------------------------%

%~~~~~~~~~~~~~~~~~~~~~~~~~~~~~~~~~~~~~~~~~~~~%
\begin{frame}{Advertisement}{Let's have an informal discussion!}
    \begin{center}
        \begin{tikzpicture}
            \node (T) {\includegraphics[width=0.95\textwidth]{AI-seminar}};
        \end{tikzpicture}
        \medskip
        
\includegraphics[width=0.5\textwidth]{Wonder}
    \end{center}
\end{frame}
%~~~~~~~~~~~~~~~~~~~~~~~~~~~~~~~~~~~~~~~~~~~~%
\AlertFrame{\Huge Exercise sheet 2}
%~~~~~~~~~~~~~~~~~~~~~~~~~~~~~~~~~~~~~~~~~~~~%
\begin{frame}[fragile]{Some guidance about the session}
    \vspace{-2mm}
    \begin{itemize}
        \item You have roughly 60 minutes\\
              \then You won't manage to do everything! \Remark{That's not expected either!}
        \item Each exercise should be solved in a dedicated script
              \PB{\Remark{Use the terminal to experiment, if needed}}
        \item Spend some time (but not more than 30 minutes) on exercise 4
        \item Exercise 7 is short and instructive: try to solve it.
        \item Exercise 2 and 3 are pedagogical about control flow and arrays\\
              \then Consider these as further practice on these topics
        \item \PS{All exercises can be solved exclusively with what we learnt so far}
        \item Note down your questions, doubts, curiosities, \ldots
        \item \alert{\textbf{Remember:} No AI of any sort for the moment!}
    \end{itemize}
    \begin{varblock}{}[0.62\textwidth]{Come up with suggestions!}
        In the closing sessions we'll solve an exercise together.
        If you have any preference, feel free to speak up.
    \end{varblock}
\end{frame}
%~~~~~~~~~~~~~~~~~~~~~~~~~~~~~~~~~~~~~~~~~~~~%
